%===========================================
% 第11章 結論（限界と今後の課題）
% 研究結果をまとめ、理論的・実証的限界を整理し、今後の研究課題（データ精緻化、動学モデル化、他産業への応用可能性）を示す。
%===========================================


%%%%%%%%%%%%%%%%%%%%%%%%%%%%%%%%%%%%%%%%%%%%%%%%%%%%%%%%%%%%%%%%%%%%%%%%%%%%%%
% 【研究の中心命題（絶対に忘れてはならない論文の背骨）】
%
% ● JEPX のスポット価格は平常時は安定している。
% ● 2021 年 1 月の高騰は「単発のショック」にすぎない。
% ● にもかかわらず、多くの新電力が倒産した。
%
% → これは「価格」では倒産を説明できないことを意味する。
%
% 【本研究の核心】
% 倒産の本質的原因は、価格ショックそのものではなく、
% 「新電力ビジネスモデルの構造的脆弱性」にある。
%
% 【構造的脆弱性の主要因】
% 1. 先渡し契約（ヘッジ）の未整備 → スポット依存 → 逆ザヤ発生
% 2. 資本力の弱さ → 数週間の逆ザヤで資金ショート
% 3. 需給調整市場コストの増加 → 小規模ほど負担が重い
% 4. 需要予測の精度不足 → 追加調達で高値掴み
% 5. JEPX の担保制度強化 → 小規模事業者は担保を積めず退出
%
% 【重要な論点】
% ● 2021 年の高騰は「引き金（トリガー）」にすぎない。
% ● 倒産の根本原因は「構造」にある。
% ● 公開データ（価格・取引量・参加者数）だけでは倒産は説明できない。
%
% 【本研究の価値（独自性）】
% ● 公開データでは見えない倒産メカニズムを解明する。
% ● 新電力の財務構造 × 市場制度 × 価格ショックの相互作用を分析する。
% ● 「倒産して当然」ではなく、「倒産する構造だった」ことを示す。
%
% 【論文全体のストーリーライン】
% 1. 市場データ分析（価格・取引量・参加者数）
%    → 表面的には安定しており、倒産は説明できない。
%
% 2. 新電力のビジネスモデル分析
%    → 構造的脆弱性を明確化。
%
% 3. 2021 年高騰の影響分析
%    → 価格ショックは「構造的弱点を露呈させた」だけ。
%
% 4. 倒産リスクの総合モデル化
%    → 財務指標 × 市場構造 × 価格ショック。
%
% 5. 政策的含意
%    → ヘッジ義務、担保制度、調整市場設計など。
%
% 【最終結論の方向性】
% 「新電力の倒産は、単発の価格ショックではなく、
% 　ビジネスモデルの構造的脆弱性が根本原因である」
%
% この方向性に論文全体を収束させること。
%%%%%%%%%%%%%%%%%%%%%%%%%%%%%%%%%%%%%%%%%%%%%%%%%%%%%%%%%%%%%%%%%%%%%%%%%%%%%%

% ------------------------------------------------------------
% 【研究の落としどころ（結論の核心）に関するメモ】
% ------------------------------------------------------------
%
% 本研究の結論は、「新電力企業の倒産は、財務指標だけでは説明できない
% “構造的リスク”によって決まる」という点にある。
%
% これは自明ではなく、また奇をてらった結論でもない。
% むしろ、第2章の倒産事例観察、第3〜5章の理論構築、
% 第6章の相互作用分析、第7章の実証分析が自然に導く帰結である。
%
% 【1. 財務指標では説明できない倒産】
% 新電力企業の倒産は、従来の倒産モデル（Altman Z、TSR/TDB、Merton、B&S）
% が前提とする「財務比率」や「収益性」だけでは説明できない。
% F-Power やウエスト電力の事例が示すように、
% 財務諸表上は急激な悪化が見えない段階でも倒産が発生している。
%
% 【2. 倒産を決定づけたのは μ・σ・S・D の“構造要因”】
% 本研究が形式化した μ（収益性）、σ（外部ショック感応度）、
% S（支援能力）、D（倒産閾値）の四要素は、
% 新電力市場の制度構造・市場構造・組織構造を反映している。
%
% 特に σ・S・D は財務諸表には現れないが、
% 倒産を決定づけた本質的な要因である。
% これらを明示的にモデル化した点が本研究の独自性である。
%
% 【3. 従来モデルには存在しない構造的メカニズム】
% Merton/B&S は株価と負債構造を前提とするが、
% 新電力企業は株価データを持たず、制度コストが外生的に変動する。
% 本研究はこの制度的特徴を μ・σ・S・D として取り込み、
% 新電力特有の倒産急増を説明する枠組みを提供する。
%
% 【4. Z_lin は“新電力専用の倒産モデル”として意味を持つ】
% 本研究の潜在倒産距離 Z_lin は、
% 新電力企業の倒産メカニズムを構造的に捉えるための指標であり、
% 従来の財務比率ベースのモデルでは捉えられない倒産リスクを説明する。
%
% 【5. 結論の位置づけ】
% 以上より、本研究の結論は次のように整理される：
%
% 「新電力企業の倒産は、財務指標ではなく、
% 市場構造・制度構造・支援構造によって決まる。
% これらを μ・σ・S・D の四要素として形式化することで、
% 従来モデルでは説明できなかった倒産急増のメカニズムを
% 理論的にも実証的にも説明できる。」
%
% この結論は自明ではなく、奇をてらったものでもなく、
% 本研究の全章が自然に導く“唯一の落としどころ”である。
%
% ------------------------------------------------------------

